\documentclass{beamer}
\usepackage{graphicx}
\usepackage{listings}
\usepackage{xcolor}
\usepackage{amsmath}
\usepackage{courier}

% Configure listings for Python code
\lstset{
    language=Python,
    basicstyle=\small\ttfamily,
    keywordstyle=\color{blue},
    stringstyle=\color{red},
    commentstyle=\color{green!60!black},
    numbers=left,
    numberstyle=\tiny,
    breaklines=true,
    frame=single,
    showstringspaces=false
}

% Better spacing for frames
\setbeamertemplate{frametitle}{\vspace{0.5em}\insertframetitle\vspace{0.5em}}
\setbeamertemplate{navigation symbols}{}
\setbeamertemplate{footline}[frame number]

\title{Introduction to Matplotlib}
\subtitle{Data Visualization in Python}
\author{Data Analysis with Python}
\date{\today}

% Enable notes
\setbeameroption{show notes}

\begin{document}

\begin{frame}
    \titlepage
\end{frame}
\note{
    Key points for introduction:
    \begin{itemize}
        \item This course covers both basic and advanced Matplotlib features
        \item Emphasize practical applications and best practices
        \item Will build from simple plots to complex visualizations
        \item Exercises will reinforce concepts through hands-on practice
    \end{itemize}
}

\begin{frame}{Course Timeline}
    \begin{itemize}
        \item 0-5 min: Introduction
        \begin{itemize}
            \item Course overview and objectives
        \end{itemize}
        \item 5-20 min: Basic Plotting
        \begin{itemize}
            \item Import and setup
            \item Simple line plots
            \item Interface approaches
        \end{itemize}
        \item 20-35 min: Subplots
        \begin{itemize}
            \item Figure/Axes hierarchy
            \item Multiple plots
            \item Layout management
        \end{itemize}
        \item 35-45 min: Plot Customization
        \item 45-60 min: Exercise 1 (Anscombe's Quartet)
        \item 60-75 min: Styles \& Export
        \item 75-85 min: Exercise 2 (Custom Styles)
        \item 85-90 min: Wrap-up and Q\&A
    \end{itemize}
\end{frame}
\note{
    Timeline management tips:
    \begin{itemize}
        \item Keep demonstrations short and focused
        \item Allow buffer time for questions
        \item Monitor exercise progress
        \item Be prepared to adjust timing based on group needs
    \end{itemize}
}

\begin{frame}[fragile]{Basic Plotting with Matplotlib}
    \begin{itemize}
        \item Matplotlib is Python's primary plotting library
        \item Two main approaches:
        \begin{itemize}
            \item \texttt{pyplot} interface (quick, simple plots)
            \item Object-oriented interface (more control)
        \end{itemize}
        \item Basic example:
    \end{itemize}
\begin{lstlisting}
import matplotlib.pyplot as plt
import numpy as np
x = np.linspace(0, 2*np.pi, 100)
plt.plot(x, np.sin(x))
plt.show()
\end{lstlisting}
\end{frame}
\note{
    Teaching tips:
    \begin{itemize}
        \item Start by explaining the difference between pyplot and object-oriented interfaces
        \item Demonstrate live coding of the simple example
        \item Show what happens if you forget plt.show()
        \item Mention common pitfalls with importing matplotlib
        \item Explain why we use numpy.linspace instead of range
    \end{itemize}
}

\begin{frame}[fragile]{Subplots}
    \begin{itemize}
        \item Use subplots for multiple plots in one figure
        \item Create using \texttt{plt.subplots(rows, cols)}
        \item Access individual plots through axes array
        \item Example:
    \end{itemize}
\begin{lstlisting}
fig, axs = plt.subplots(2, 1)
axs[0].plot(x, np.sin(x))
axs[1].plot(x, np.cos(x))
\end{lstlisting}
\end{frame}
\note{
    Important concepts:
    \begin{itemize}
        \item Explain figure vs axes hierarchy
        \item Show how to access individual subplots using array indexing
        \item Demonstrate different subplot layouts (2x2, 3x1, etc.)
        \item Mention tight\_layout() for fixing overlapping issues
        \item Show how to share axes between subplots
    \end{itemize}
}

\begin{frame}[fragile]{Plot Customization}
    \begin{itemize}
        \item Labels and titles
        \item Colors and styles
        \item Markers and line properties
        \item Legend and grid
        \item Example:
    \end{itemize}

\begin{lstlisting}
ax.plot(x, y, color='blue', label='data')
ax.set_xlabel('X axis')
ax.set_ylabel('Y axis')
ax.set_title('My Plot')
ax.legend()
ax.grid(True)
\end{lstlisting}
\end{frame}
\note{
    Customization tips:
    \begin{itemize}
        \item Show real-world examples of when to use different styles
        \item Demonstrate interactive customization in Jupyter
        \item Discuss color choices for accessibility
        \item Show how to save custom styles for reuse
        \item Mention journal-specific formatting requirements
    \end{itemize}
}

\begin{frame}{Exercise 1: Anscombe's Quartet}
    \begin{itemize}
        \item Famous dataset showing importance of visualization
        \item Four datasets with identical statistics
        \item Different visual patterns
        \item Tasks:
        \begin{itemize}
            \item Load data from JSON
            \item Create scatter plots
            \item Add regression lines
            \item Compare results
        \end{itemize}
    \end{itemize}
\end{frame}
\note{
    Exercise guidance:
    \begin{itemize}
        \item Let students attempt loading JSON data first
        \item Point out that all datasets have identical:
        \begin{itemize}
            \item Mean x $\approx$ 9.0
            \item Mean y $\approx$ 7.5
            \item Correlation coefficient $\approx$ 0.816
            \item Linear regression: y = 3 + 0.5x
        \end{itemize}
        \item Discuss why visualization is crucial for data analysis
    \end{itemize}
}

\begin{frame}[fragile]{Exercise 2: Custom Plotting Styles}
    \begin{itemize}
        \item Create professional-looking plots
        \item Matplotlib style sheets
        \item Components to customize:
        \begin{itemize}
            \item Figure size and DPI
            \item Font sizes and families
            \item Line styles and colors
            \item Grid properties
        \end{itemize}
    \end{itemize}
    \vspace{1em}
    \begin{lstlisting}
plt.style.use('seaborn')  # Example style
    \end{lstlisting}
\end{frame}
\note{
    Style implementation tips:
    \begin{itemize}
        \item Show examples of built-in styles first
        \item Demonstrate creating a custom style file
        \item Discuss common style elements for scientific plots
        \item Show how to combine multiple style sheets
        \item Reference matplotlib style sheet reference online
    \end{itemize}
}

\begin{frame}[fragile]{Saving Plots}
    \begin{itemize}
        \item Use \texttt{plt.savefig()} or \texttt{fig.savefig()}
        \item Common formats:
        \begin{itemize}
            \item PNG (raster) -- Best for web
            \item SVG (vector) -- Scalable graphics
            \item PDF (vector) -- Publication quality
        \end{itemize}
        \item Control DPI and size:
        \begin{lstlisting}
fig.savefig('plot.svg',
            dpi=300,
            bbox_inches='tight',
            format='svg')
        \end{lstlisting}
    \end{itemize}
\end{frame}
\note{
    Export guidelines:
    \begin{itemize}
        \item Explain when to use each file format:
        \begin{itemize}
            \item PNG: web, presentations
            \item SVG: web with zoom, figures needing editing
            \item PDF: publication submissions
        \end{itemize}
        \item Discuss resolution requirements for different purposes
        \item Show how to handle font embedding
        \item Mention size limits for journal submissions
    \end{itemize}
}

\end{document}
